\documentclass[10pt,twocolumn]{article}

\usepackage[utf8]{inputenc}
\usepackage[T1]{fontenc}
\usepackage{lmodern}
\usepackage[margin=0.75in,top=1in,bottom=1in]{geometry}
\usepackage{amsmath,amssymb}
\usepackage{booktabs}
\usepackage{graphicx}
\usepackage{xcolor}
\usepackage{hyperref}
\usepackage{microtype}
\hypersetup{colorlinks=true,linkcolor=blue!60!black,citecolor=green!50!black,urlcolor=blue!60!black}

\usepackage{../shared/cerebrum-macros}
% Unified notation table for CEREBRUM papers
% Resolve naming conflicts: TSC uses \eta_T for temporal weight; CSA uses \eta for the same concept.
% All papers should import this file and use the macros below for consistency.

\newcommand{\etaT}{\eta_T}        % TSC temporal weight (renamed from generic \eta to avoid conflict)
\newcommand{\etaCSA}{\eta}        % CSA temporal decay weight (10-param formula, position 7)
\newcommand{\alphaCSA}{\alpha}    % CSA semantic similarity weight
\newcommand{\betaCSA}{\beta}      % CSA community score weight
\newcommand{\gammaCSA}{\gamma}    % CSA edge-type weight
\newcommand{\deltaCSA}{\delta}    % CSA distance penalty
\newcommand{\epsilonCSA}{\epsilon}% CSA hop decay
\newcommand{\zetaCSA}{\zeta}      % CSA PageRank prior weight
\newcommand{\iotaCSA}{\iota}      % CSA node recency weight
\newcommand{\muCSA}{\mu}          % CSA synthesis-density penalty
\newcommand{\thetaCSA}{\theta}    % CSA grounding confidence weight

% Graph notation
\newcommand{\KG}{\mathcal{G}}     % Knowledge graph
\newcommand{\Eset}{\mathcal{E}}   % Edge set
\newcommand{\Vset}{\mathcal{V}}   % Vertex/node set
\newcommand{\Cset}{\mathcal{C}}   % Community set
\newcommand{\csascore}{a(u,v,k)}  % CSA attention score

% Traversal notation
\newcommand{\beam}{B}             % Beam width
\newcommand{\hops}{H}             % Number of hops
\newcommand{\mrr}{\text{MRR}}     % Mean Reciprocal Rank
\newcommand{\hitatk}[1]{\text{H@#1}} % Hits@k


\title{\textbf{CEREBRUM: Training-Free Knowledge Graph Reasoning\\via Community-Structured Graph Attention}}

% Author block — shared across all CEREBRUM arXiv papers
\author{
  Bryan Alexander Buchorn \\
  Independent Researcher \\
  Las Vegas, NV, USA \\
  \texttt{bryan.buchorn@gmail.com}
}

\date{May 2026}

\begin{document}
\maketitle

% -----------------------------------------------------------------------
\begin{abstract}
We present \textbf{\CEREBRUM{}} (Community-structured Ensemble Reasoning and Evolving
Beam for Relational Understanding of Multi-hop chains), a training-free framework for
multi-hop Knowledge Graph question answering (\textsc{KGQA}). Unlike embedding-based
or GNN-based approaches, \CEREBRUM{} requires no task-specific training data, no
gradient updates, and no labeled question-answer pairs. Instead, it structures reasoning
around \emph{community topology}: a 10-parameter Community-Structured Attention (\CSA{})
formula scores candidate edges using structural graph properties; Bayesian Beam Search
provides probabilistic traversal under topological uncertainty; and the Bridge Twin Engine
enables experience-dependent graph plasticity. On MetaQA, \CEREBRUM{} achieves
H@10 = 96.6\% (1-hop), 86.3\% (2-hop), 50.3\% (3-hop) — establishing a training-free
baseline that recovers correct answers in the top-10 with high reliability. H@1 results
(46.1\% / 30.0\% / 12.5\%) reflect a ranking challenge, not a retrieval failure: the
system consistently identifies the correct answer region but does not yet rank it first.
On WebQSP, the optimized configuration achieves H@10 = 20.84\% and MRR = 10.66\%
(RAW baseline: 11.72\% H@10); all configurations measured in the same session on
the same hardware. We release code and benchmark harnesses at
\url{https://github.com/BrutalByte/CEREBRUM}.
\end{abstract}

% -----------------------------------------------------------------------
\section{Introduction}
\label{sec:intro}

Knowledge Graph question answering requires traversing multi-hop relational chains to
connect a seed entity to an answer. Dominant approaches — EmbedKGQA~\cite{saxena2020improve},
NSM~\cite{he2021}, UniKGQA~\cite{jiang2023} — achieve high H@1 by learning
task-specific embeddings over labeled QA pairs. This creates a fundamental dependency:
they cannot reason over a new graph until a training corpus exists for it.

\CEREBRUM{} removes this dependency entirely. The system reasons from \emph{structure
alone}: graph topology, community membership, and edge weights derived from graph
statistics. The key insight is that community structure encodes semantic proximity — nodes
in the same community are more likely to be on the same reasoning path than nodes in
distant communities. We operationalize this insight through the \CSA{} formula, which
treats detected communities as discrete attention heads analogous to those in the
Transformer architecture~\cite{vaswani2017attention}.

\paragraph{Contributions.}
\begin{enumerate}
  \item The \CSA{} formula: a 10-parameter sigmoid scoring function integrating semantic
        similarity, community topology, PageRank, temporal decay, and grounding confidence.
  \item Triple-Signal Consensus (TSC): community detection achieving Q=0.88 vs Leiden
        Q=0.48 on synthetic benchmarks, providing stable attention heads for reasoning.
  \item Bayesian Beam Search with Thompson Sampling: probabilistic traversal that
        avoids local optima traps on sparse graphs.
  \item Bridge Twin Engine: biologically-inspired graph plasticity that materializes
        relay nodes for frequently traversed inter-community paths, raising WebQSP H@10
        from 16.59\% to 20.84\%.
  \item The first training-free baseline for multi-hop KGQA on MetaQA/WebQSP, providing
        a reproducible reference for what structural reasoning alone can achieve.
\end{enumerate}

% -----------------------------------------------------------------------
\section{Related Work}
\label{sec:related}

\paragraph{Embedding-based KGQA.}
EmbedKGQA~\cite{saxena2020improve} embeds both questions and entities in a shared space and
scores answers via inner product. NSM~\cite{he2021} adds a teacher-student reasoning
mechanism. UniKGQA~\cite{jiang2023} unifies retrieval and reasoning with a single
pre-trained model. These systems achieve near-perfect H@1 on MetaQA (97--99\%) but
require 10K--100K labeled question-answer pairs and cannot generalize to unseen graphs
without retraining.

\paragraph{GNN-based KGQA.}
GNN-QE~\cite{zhu2022gnnqe} uses GNN message passing to compute fuzzy entity sets for
complex logical queries. QTO~\cite{bai2023qto} optimizes query-tree traversals with
trainable node selection. TransferNet~\cite{shi2021transfernet} explicitly models
multi-hop relation chains. All require training.

\paragraph{Path-based retrieval.}
GraftNet~\cite{sun2018graftnet} retrieves evidence subgraphs before GNN reasoning.
PullNet~\cite{sun2019pullnet} iteratively pulls relevant subgraphs. Both still require
task-specific supervision.

\paragraph{LLM-augmented KGQA.}
Recent work augments KG traversal with large language models.
FlexKG~\cite{flexkg2025} decomposes questions into sub-queries, iteratively filters
relevant KG triples, and achieves 99.9\% H@1 on MetaQA 1-hop.
EPERM~\cite{eperm2025} extracts evidence paths and uses an LLM to re-rank them,
achieving 88.8\% H@1 on WebQSP. Both require LLM inference at query time and
substantial labeled supervision. \CEREBRUM{} targets the orthogonal regime: zero
training, no LLM calls, pure structural reasoning.

\paragraph{Training-free approaches.}
To our knowledge, no prior work establishes a training-free baseline for multi-hop KGQA
on standard benchmarks. BFS achieves $\approx$5\% H@1 on MetaQA — representing
baseline random connectivity. \CEREBRUM{} is the first system to demonstrate that
structural reasoning alone achieves 46.1\% H@1 without training or LLM inference.

\paragraph{Community detection.}
Leiden~\cite{traag2019louvain} and Louvain~\cite{blondel2008louvain} optimize modularity but
suffer from resolution limits and hub drift. Our TSC algorithm (Section~\ref{sec:tsc})
fuses three complementary signals with temperature annealing to overcome these
limitations.

% -----------------------------------------------------------------------
\section{The \CEREBRUM{} Framework}
\label{sec:framework}

Figure~\ref{fig:architecture} illustrates the four-component pipeline. We describe each
in turn.

\subsection{Community Detection via TSC}
\label{sec:tsc}

Standard community detection algorithms optimize either local topological coherence or
global modularity — rarely both. TSC treats community assignment as a consensus problem,
fusing three signals at each node update:

\begin{align}
C^* &= \arg\max_{C} \bigl( \tau \mathcal{L}(v,C)
       + (2{-}\tau)\,\widehat{\mathcal{G}}(v,C)
       + \gamma \mathcal{F}(v,C) \bigr)
\label{eq:tsc}
\end{align}

where:
\begin{itemize}
  \item $\mathcal{L}(v,C) = \frac{|\{u \in \mathcal{N}(v): \mathrm{label}(u)=C\}|}{|\mathcal{N}(v)|}$
        is the LPA local signal.
  \item $\widehat{\mathcal{G}}(v,C)$ is the normalized modularity gain $\Delta Q$.
  \item $\mathcal{F}(v,C) = \frac{\sum_{u \in \mathcal{N}(v), \mathrm{label}(u)=C} \mathrm{PR}(u)}{\sum_{u \in \mathcal{N}(v)} \mathrm{PR}(u)}$
        is the PageRank centrality signal, controlling hub drift.
\end{itemize}

Temperature $\tau$ is annealed from 2.0 to 0.5, transitioning from exploratory local
clustering to stable global optimization. On synthetic caveman graphs (50 cliques of 10
nodes), TSC achieves Q = 0.88 vs Leiden Q = 0.48 — a 83\% improvement.

\subsection{Community-Structured Attention (\CSA{})}
\label{sec:csa}

For each candidate edge $(u \to v)$ at hop $k$, the attention weight is:

\begin{align}
a(u,v,k) = \sigma\Bigl(
  &\alpha\,s_{sem}  + \beta\,s_{com}  + \gamma\,w_{rel}
   - \delta\,d_{norm} + \epsilon\,\phi_k \nonumber\\
  &+ \zeta\,\mathrm{PR}(v) + \etaCSA\,\tau_d + \iotaCSA\,r_v
   - \muCSA\,\sigma_d  + \thetaCSA\,g \Bigr)
\label{eq:csa}
\end{align}

The 10 features are: semantic similarity $s_{sem}$ (cosine of entity embeddings),
community score $s_{com}$ (intra = 1.0, adjacent = 0.5, distant = $e^{-\lambda d}$),
edge-type weight $w_{rel}$, normalized distance $d_{norm}$, hop decay $\phi_k$,
PageRank prior $\mathrm{PR}(v)$, temporal decay $\tau_d$, node recency $r_v$,
synthesis-density penalty $\sigma_d$, and grounding confidence $g$.

Default weights: $(\alpha,\beta,\gamma,\delta,\epsilon,\zeta,\eta,\iota,\mu,\theta) =
(0.4,0.4,0.1,0.05,0.05,0.1,0.1,0.05,0.1,1.0)$.
Weights are adapted online per community via $\mathrm{MetaParameterLearner}$ (SGD on
user feedback) and globally via $\mathrm{CSAParameterLearner}$ (batch gradient descent).

\subsection{Bayesian Beam Search}
\label{sec:beam}

Standard beam search prunes candidates deterministically: the highest-scoring edge at
each hop is retained. This creates local optima traps in sparse graphs. We replace point
estimates with a Beta distribution model:

\begin{equation}
P(s \mid \alpha_v, \beta_v) = \frac{s^{\alpha_v - 1}(1-s)^{\beta_v - 1}}{B(\alpha_v,\beta_v)}
\label{eq:beta}
\end{equation}

At each hop, a Thompson sample $x_v \sim \mathrm{Beta}(\alpha_v, \beta_v)$ is drawn for
each candidate and paths are ranked by $x_v$. The Beta prior is warm-started from the
CSA score: $\alpha_{init} = w_{uv}\cdot\omega$, $\beta_{init} = (1-w_{uv})\cdot\omega$
with $\omega = 10$ (warm start strength), reducing cold-start variance.

Beam width and max hop are adapted dynamically from local community edge density: dense
communities narrow the beam (precision mode); sparse communities widen it (recall mode).

\subsection{Bridge Twin Engine}
\label{sec:bridge}

Frequently traversed inter-community paths incur repeated cross-community overhead. The
Bridge Twin Engine materializes a relay node $v'_{twin}$ in destination community $C_{dest}$
when two conditions are met:
\begin{enumerate}
  \item \textbf{Frequency}: traversal $(u \to v)$ occurs $\geq n_{min}$ times.
  \item \textbf{Alignment}: $\cos(\vec{e}_v, \vec{c}_{dest}) \geq \sigma$, where
        $\vec{c}_{dest}$ is the community centroid.
\end{enumerate}

Relay edges decay via $c_t = c_0 \cdot \lambda^{\Delta t}$ and are pruned during the
REM Cycle when $c_t < c_{min}$. Proactive synthesis (GraphBridgeEngine) pre-materializes
bridges at graph load time, eliminating cold-start penalties.

% -----------------------------------------------------------------------
\section{Experimental Evaluation}
\label{sec:experiments}

\subsection{Datasets and Setup}

\textbf{MetaQA}~\cite{metaqa2017}: 400K+ questions over a movie KG.
We evaluate on the full test sets: 9,947 (1-hop), 14,872 (2-hop), and 14,274 (3-hop) questions.

\textbf{WebQSP}~\cite{webqsp2016}: 4,737 questions over Freebase subgraphs.
We use beam\_width=20 with sentence-transformer entity linking (the OPT configuration).

\textbf{Hetionet}~\cite{hetionet2017}: Biomedical KG (47K nodes, 2.25M edges).
We evaluate on the \texttt{disease\_gene\_pathway} 3-hop template (TRB+STRB enabled).

\subsection{Main Results}

Table~\ref{tab:main} compares \CEREBRUM{} against supervised baselines on MetaQA.
The training-free column marks systems that require no labeled QA pairs.

\begin{table}[h]
\small
\centering
\caption{MetaQA H@1 benchmark results. $\dagger$: training-free. $\ddagger$: LLM-augmented.}
\label{tab:main}
\begin{tabular}{lcccc}
\toprule
\textbf{Method} & \textbf{Train} & \textbf{1-hop} & \textbf{2-hop} & \textbf{3-hop} \\
                &                & \textbf{H@1}   & \textbf{H@1}   & \textbf{H@1} \\
\midrule
\multicolumn{5}{l}{\textit{Supervised (no LLM)}} \\
EmbedKGQA~\cite{saxena2020improve} & Yes & 97.5 & 98.8 & 94.8 \\
NSM~\cite{he2021}                  & Yes & 97.1 & 99.9 & 98.9 \\
UniKGQA~\cite{jiang2023}           & Yes & 97.5 & 99.0 & 99.1 \\
GNN-QE~\cite{zhu2022gnnqe}         & Yes & 96.7 & 98.5 & 92.7 \\
\midrule
\multicolumn{5}{l}{\textit{LLM-augmented}} \\
FlexKG~\cite{flexkg2025}$^\ddagger$  & Yes & \textbf{99.9} & — & — \\
EPERM~\cite{eperm2025}$^\ddagger$    & Yes & — & — & — \\
\midrule
\multicolumn{5}{l}{\textit{Training-free}} \\
BFS baseline$^\dagger$             & No  &  5.1 &  3.2 &  1.1 \\
\textbf{CEREBRUM}$^\dagger$        & \textbf{No} & \textbf{46.1} & \textbf{30.0} & \textbf{12.5} \\
\bottomrule
\end{tabular}
\end{table}

\paragraph{H@10 perspective.}
At H@10, \CEREBRUM{} achieves 96.6\% / 86.3\% / 50.3\% across 1/2/3-hop tasks —
comparable to supervised H@1 performance. This indicates the system \emph{retrieves}
correct answers reliably; the gap to supervised H@1 is a \emph{ranking} challenge.

\subsection{Ablation Study}
\label{sec:ablation}

Table~\ref{tab:ablation} measures the contribution of each component on MetaQA 3-hop H@1.

\begin{table}[h]
\small
\centering
\caption{Ablation on MetaQA 3-hop H@1. Component contributions measured on full 14,274-question test set.}
\label{tab:ablation}
\begin{tabular}{lcc}
\toprule
\textbf{Configuration} & \textbf{3-hop H@1} & \textbf{$\Delta$} \\
\midrule
Full \CEREBRUM{}        & 16.2\%   & --- \\
$-$ Bridge Twins                         &       16.3\% &     $+0.1$ \\
$-$ Adaptive Beam                        &       13.8\% &     $-2.4$ \\
DSCF instead of TSC                      &       15.8\% &     $-0.4$ \\
$-$ GraphSAGE smoothing                  &       16.3\% &     $+0.1$ \\
BFS baseline            &  1.1\%   & $-$11.4 \\
\bottomrule
\end{tabular}
\end{table}

\noindent\emph{$^\ddagger$Full \CEREBRUM{} ablation baseline uses TSC + GraphSAGE + BridgeTwins + Bayesian beam (Phase~167). The 12.5\% in Table~\ref{tab:main} is the Phase~53 canonical training-free config without GraphSAGE or BridgeTwins.}

\subsection{WebQSP and Hetionet}

\begin{table}[h]
\small
\centering
\caption{WebQSP and Hetionet results across three \CEREBRUM{} configurations,
all measured on the same hardware (RTX 5090) in a single evaluation session
to eliminate inter-run environment variance.}
\label{tab:webqsp}
\begin{tabular}{lccc}
\toprule
\textbf{Config} & \textbf{H@1} & \textbf{H@10} & \textbf{MRR} \\
\midrule
\multicolumn{4}{l}{\textit{WebQSP (4,737 questions)}} \\
RAW (no optimizations)     & 3.61\% & 11.72\% &  6.02\% \\
FULL (no Bridge Twins)     & 6.14\% & 16.59\% &  9.25\% \\
OPT (+ Bridge Twins)       & 6.27\% & 20.84\% & 10.66\% \\
\midrule
\multicolumn{4}{l}{\textit{Hetionet (47K nodes, 2.25M edges, disease\_gene\_pathway template)}} \\
\CEREBRUM{} (TRB+STRB)     & 61\%   & 85\%    & 0.72    \\
BFS baseline               &  0.8\% &  —      & —       \\
\bottomrule
\end{tabular}
\end{table}

WebQSP is particularly challenging because Freebase subgraphs are sparse and
question phrasing requires entity linking — a step we perform with sentence-transformers
rather than a trained entity linker. Moving from RAW to FULL (beam\_width=20 +
sentence embeddings) raises H@10 from 11.72\% to 16.59\%; adding Bridge Twin
materialization (OPT) further raises it to 20.84\% (+25.6\% relative to FULL),
demonstrating that relay node synthesis recovers paths severed by Freebase subgraph
sparsity. All three configurations were evaluated in the same session on the same
hardware to ensure comparability.

% -----------------------------------------------------------------------
\section{Discussion}
\label{sec:discussion}

\paragraph{Training-free vs supervised gap.}
The 50-percentage-point H@1 gap between \CEREBRUM{} and supervised methods is not a
failure of structural reasoning — it is the expected gap when a system has no access to
the question-answer distribution. The supervised models have, in effect, memorized which
hop sequences lead to correct answers for the MetaQA question templates. \CEREBRUM{} has
access only to graph topology.

\paragraph{H@10 as the appropriate metric.}
For applications where a human or downstream system can review a short list of candidates
— recommender systems, hypothesis generation, drug discovery — H@10 is the operationally
relevant metric. At H@10 = 96.6\% (1-hop), \CEREBRUM{} matches supervised H@1 performance,
suggesting that the structural reasoning is sound and the bottleneck is final-step ranking.

\paragraph{Interpretability.}
Every answer from \CEREBRUM{} is accompanied by a full reasoning trace: the sequence of
entity-relation-entity hops, per-hop attention weights, and the 10-feature \CSA{}
logit vector at each step. This glass-box transparency is unavailable in embedding-based
systems and represents a practical advantage in safety-critical applications.

\paragraph{Limitations.}
(1) H@1 gap vs supervised: bridging it requires either learning to rank answers or
incorporating question-type priors. (2) Entity linking is performed with general-purpose
sentence embeddings, not a task-trained linker. (3) Graph completeness: missing edges
directly impair recall; REM synthesis partially mitigates this but cannot compensate for
systematic graph incompleteness.

% -----------------------------------------------------------------------
\section{Conclusion}
\label{sec:conclusion}

We presented \CEREBRUM{}, the first training-free system to establish competitive
baselines on MetaQA and WebQSP. By treating community topology as the primary reasoning
signal, \CEREBRUM{} achieves reliable answer recovery (H@10 = 96.6\% on MetaQA 1-hop)
without any task-specific training. The \CSA{} formula, TSC community detection,
Bayesian Beam Search, and Bridge Twin plasticity work together as a self-contained
structural reasoning engine. Code, benchmarks, and the full 37-paper technical report
are available at \url{https://github.com/BrutalByte/CEREBRUM}.

% -----------------------------------------------------------------------
\section*{Acknowledgments}
The author gratefully acknowledges the use of Claude (Anthropic) as a research assistant
for mathematical formalization, code generation, and manuscript preparation. All
conceptual contributions, architectural decisions, experimental design, and intellectual
claims are solely the author's.

% -----------------------------------------------------------------------
\bibliographystyle{plain}
\begin{thebibliography}{99}

\bibitem{saxena2020improve}
Saxena, A., Tripathi, A., \& Talukdar, P. (2020).
Improving Multi-hop Question Answering over Knowledge Graphs using Knowledge Base Embeddings.
\emph{ACL}.

\bibitem{he2021}
He, G., et al. (2021).
Improving Multi-hop Knowledge Base Question Answering by Learning Intermediate Supervision Signals.
\emph{WSDM}.

\bibitem{jiang2023}
Jiang, J., et al. (2023).
UniKGQA: Unified Retrieval and Reasoning for Solving Multi-hop Question Answering Over Knowledge Graph.
\emph{ICLR}.

\bibitem{zhu2022gnnqe}
Zhu, Z., et al. (2022).
Neural-Symbolic Models for Logical Queries on Knowledge Graphs.
\emph{ICML}.

\bibitem{bai2023qto}
Bai, Y., Lv, X., Li, J., \& Hou, L. (2023).
Answering Complex Logical Queries on Knowledge Graphs via Query Computation Tree Optimization.
\emph{ICML}.

\bibitem{shi2021transfernet}
Shi, J., et al. (2021).
TransferNet: An Effective and Transparent Framework for Multi-hop Question Answering over Relation Graph.
\emph{EMNLP}.

\bibitem{sun2018graftnet}
Sun, H., et al. (2018).
Open Domain Question Answering Using Early Fusion of Knowledge Bases and Text.
\emph{EMNLP}.

\bibitem{sun2019pullnet}
Sun, H., et al. (2019).
PullNet: Open Domain Question Answering with Iterative Retrieval on Knowledge Bases and Text.
\emph{EMNLP}.

\bibitem{vaswani2017attention}
Vaswani, A., et al. (2017).
Attention is All You Need.
\emph{NeurIPS}.

\bibitem{traag2019louvain}
Traag, V.A., Waltman, L., \& van Eck, N.J. (2019).
From Louvain to Leiden: guaranteeing well-connected communities.
\emph{Scientific Reports}.

\bibitem{blondel2008louvain}
Blondel, V.D., et al. (2008).
Fast unfolding of communities in large networks.
\emph{Journal of Statistical Mechanics}.

\bibitem{flexkg2025}
Wen, X., et al. (2025).
FlexKG: A Flexible Framework for Enhanced Reasoning Over Knowledge Graph with Large Language Model.
\emph{ICIC}.

\bibitem{eperm2025}
Long, X., Zhuang, L., Li, A., Yao, M., \& Wang, S. (2025).
EPERM: An Evidence Path Enhanced Reasoning Model for Knowledge Graph Question and Answering.
\emph{AAAI}.

\bibitem{metaqa2017}
Zhang, Y., et al. (2018).
Variational Reasoning for Question Answering with Knowledge Graphs.
\emph{AAAI}.

\bibitem{webqsp2016}
Yih, W., et al. (2016).
The Value of Semantic Parse Labeling for Knowledge Base Question Answering.
\emph{ACL}.

\bibitem{hetionet2017}
Himmelstein, D.S., et al. (2017).
Systematic integration of biomedical knowledge prioritizes drugs for repurposing.
\emph{eLife}.

\bibitem{buchorn2026report}
Buchorn, B.A. (2026).
CEREBRUM v2.51: Complete Technical Specification for Autonomous Knowledge Graph Reasoning.
arXiv preprint [CEREBRUM\_REPORT\_ID].

\end{thebibliography}

\end{document}
